% Part: computability
% Chapter: recursive-functions
% Section: trees

\documentclass[../../../include/open-logic-section]{subfiles}

\begin{document}

\olfileid{cmp}{rec}{tre}
\olsection{درخت}

کدی کدی درختاں نوں قدرتی عدداں دے روپ وچ ظاہر کرنا
مفید ہوندا اے؛ جداں اسی سلسلیاں نوں عدداں دے روپ وچ،
تے اوہناں دیاں خاصیتاں تے عملاں نوں اوہناں دے کوڈاں
اُتے بنیادی بازگشتی تعلقاں تے فنکشناں دے روپ وچ ظاہر
کر سکدے آں۔ ایتھے اسی سلسلے تے اوہناں دے کوڈ ورتاں گے۔
درخت یا تاں اکّلی گرہ ہو سکدا اے (ممکن اے اوہدے
نال کوئی لیبل ہووے)، یا اک گرہ جیہدے نال ذیلی درخت
جُڑے ہون (اوہدے نال وی لیبل ہو سکدا اے)۔
ایس گرہ نوں درخت دی \emph{جڑ} کہندے آں، تے اوہدے نال
سِدھے جُڑے ذیلی درختاں نوں \emph{براہِ راست ذیلی درخت}۔

اسی درختاں نوں بازگشتی طور اُتے سلسلہ
$\tuple{k, d_1, \dots, d_k}$ نال کوڈ کردے آں،
جتھے $k$ براہِ راست ذیلی درختاں دی تعداد اے،
تے $d_1$، \dots،~$d_k$ اوہناں دے کوڈ نیں۔
جے گرہواں دے لیبل وی ہون، تاں اوہناں نوں براہِ راست
ذیلی درختاں توں بعد شامل کیتا جا سکدا اے۔ ایس لئی
صرف اک گرہ تے لیبل~$l$ والا درخت
$\tuple{0,l}$ نال کوڈ ہووے گا؛ تے اک جڑ (لیبل~$l_1$)
نال جُڑیاں دو اکّلی گرہواں (لیبل $l_2$، $l_3$)
والا درخت $\tuple{2,
  \tuple{0, l_2}, \tuple{0, l_3}, l_1}$ نال کوڈ ہووے گا۔

\begin{prop}
  \ollabel{prop:subtreeseq}
فنکشن $\fn{SubtreeSeq}(t)$، جیہڑا اک سلسلے دا کوڈ
واپس دیندا اے جیہدے رکن کوڈ~$t$ والے درخت دے سارے
ذیلی درختاں دے کوڈ ہون، بنیادی بازگشتی اے۔
\end{prop}

\begin{proof}
پہلاں غور کرو کہ $\fn{ISubtrees}(t) = \fn{subseq}(t, 1, (t)_0)$
بنیادی بازگشتی اے تے درخت~$t$ دے براہِ راست
ذیلی درختاں دے کوڈ واپس دیندا اے۔ ہُن اسی اک مددگار
فنکشن $\fn{hSubtreeSeq}(t,n)$ متعین کر سکدے آں
جیہڑا جڑ توں $n$ قدم تک دے سارے ذیلی درختاں دا
سلسلہ حساب کردا اے۔ درخت~$t$ دی جڑ توں $0$
قدم تک دے ذیلی درختاں دا سلسلہ—یعنی اوہ جیہڑا
درخت~$t$ دی جڑ توں شروع ہوندا اے—صرف~$t$
اُتے مشتمل اے۔ سطح~$n+1$ تک درخت $t$ دے سارے
ذیلی درختاں دا سلسلہ لین لئی، سطح $n$ تک جمع
ذیلی درختاں نال سطح $n$ تک دے انہاں ساریاں دے
براہِ راست ذیلی درختاں دا سلسلہ جوڑ دیندے آں۔
ایہہ سارے اکٹھے کرن لئی غور کرو کہ جے $f(x)$
بنیادی بازگشتی فنکشن اے جیہڑا سلسلیاں دے کوڈ واپس
دیندا اے، تاں $g_f(s, k) = f((s)_0) \concat
\dots \concat f((s)_k)$ وی بنیادی بازگشتی اے:
    \begin{align*}
      g(s, 0) & = f((s)_0)\\
      g(s, k+1) & = g(s, k) \concat f((s)_{k+1})
    \end{align*}
    مثال لئی، جے $s$ درختاں دا سلسلہ اے، تاں
    $h(s) = g_{\fn{ISubtrees}}(s, \len{s})$ سلسلے~$s$
    دے رکناں دے براہِ راست ذیلی درختاں دا سلسلہ دیندا
    اے۔ اسی ایہہ ورت کے~$\fn{hSubtreeSeq}$ نوں
    ایویں متعین کر سکدے آں:
    \begin{align*}
      \fn{hSubtreeSeq}(t, 0) & = \tuple{t} \\
      \fn{hSubtreeSeq}(t, n+1) & = \fn{hSubtreeSeq}(t, n) \concat
      h(\fn{hSubtreeSeq}(t, n)).
    \end{align*}
    کوڈ~$t$ والے درخت وچ ذیلی درختاں دی زیادہ توں
    زیادہ سطح، یعنی جڑ تے کسے پتے والی گرہ دے
    وچکار زیادہ توں زیادہ فاصلہ، کوڈ~$t$ توں ودھ نہیں۔
    ایس لئی کوڈ~$t$ والے درخت دے سارے ذیلی درختاں
    دے کوڈاں دا اک سلسلہ $\fn{hSubtreeSeq}(t, t)$
    نال مل جاندا اے۔
\end{proof}

\begin{prob}
  ثبوت وچ $\fn{hSubtreeSeq}$ دی تعریف،
  \olref[cmp][rec][tre]{prop:subtreeseq} وچ، عام طور
  اُتے دہرائیاں شامل کردی اے۔ اک وکھری تعریف دیو
  جیہڑی یقینی بنائے کہ نتیجے والی فہرست وچ کسے
  ذیلی درخت دا کوڈ صرف اک وار آوے۔
\end{prob}

\end{document}
